\section{Related Work}
Natural-reading eye-tracking corpora support analyses of reading behavior under
connected text rather than isolated word presentation \cite{kennedy2003d,kliegl2006tracking}.
CopCo provides Danish natural-reading material that can be aligned with gaze,
linguistic, and language-model features. We treat the current labels as operational
research labels and do not claim clinical diagnosis.

Reader-level prediction from eye movements has been explored in dyslexia and reading
difficulty settings, including Danish work on gaze-based prediction. Our contribution is
not a broader model search, but an interpretable confirmatory package that separates
text exposure from participant-level predictability sensitivity. LM surprisal has long
been connected to reading time and processing difficulty \cite{hale2001probabilistic,
levy2008expectation,smith2013effect}. We use DFM surprisal and entropy as contextual
predictability features, then ask whether participant gaze costs vary with these
features. Danish vocalic and boundary-opacity literature motivates the secondary
orthographic boundary-opacity analysis, but the current labels are orthographic proxies
rather than pronunciation-aware phonological labels. Explainable reader-level prediction
requires grouped validation, leakage controls, and feature-stability analysis rather
than random word-level splits.
